\section{Conclusion}\label{sec:conclusion}

We asked what pacing minimizes 200-yard freestyle time under realistic
energetic and fatigue constraints, and whether competitive swimmers use it.
On the first question the answer has a structure that does not depend on
parameters: under a fixed energy budget with a velocity-only cost, even
pacing is the unique optimum for every cost exponent above one, and with
position-coupled economy decay the optimal split allocation
$t_i\propto w_i^{1/p}$ is invariant to aerobic capacity, anaerobic
capacity, drag and efficiency. Fitness sets the speed of the race and not
its shape, oxygen kinetics slow the race and do not bend it, and every
departure of a real race from even pacing is evidence about a mechanism
outside the budget.

On the second question we built a verified pipeline from official result
sheets to anonymized, checksum-registered data, wrote and froze an analysis
plan before any race was in hand, and tested five fatigue mechanisms on 345
races by male swimmers aged 15--18 with a swimmer-grouped held-out set
opened once. The races are positively split; the even and negative-split
mechanisms are rejected on held-out accuracy at the elite-anchored start
credit; the two positive-split mechanisms, position-coupled cost and a
fatigue-coupled velocity ceiling, tie, because the observed fade is
front-loaded in kind but between them in degree. The ranking among the
surviving models changes inside the literature-anchored band of dive-start
credits, which under the registered plan means the data cannot distinguish
them given start uncertainty; the hypothesis that swimmers nearer the
predicted optimum perform better relative to expectation returned the null
the plan anticipated; and the mean profile carries a finishing kick that no
monotone mechanism can produce.

The next steps are set by those findings rather than by ambition. The
dive-start credit must be measured, with 15~m splits, or at least modelled
as a function of pace, before any fatigue coupling fitted to split sheets
can be interpreted. The finishing kick calls for a model with uncertainty
in the swimmer's state, which is a different and more interesting
optimization problem than the deterministic one solved here. And the
deviation--performance question needs many races of few swimmers with a
real expectation model, not few races of many swimmers with a stale one.
The code, the analysis plan with its amendment log, the dataset register
with the digests of every official source, and every number in this paper
are in the project repository, so that the next attempt can start where
this one stopped rather than where it started.
